
\section{Discussion}\label{sec:discussion}
The engineering-informatics contribution assigns distinct computational and evidential consequences to source lineage, declared profiles, processing service and statistical events. In the secondary request analysis, changing a trajectory requirement to an observation requirement changes the selected geometry policy, recovers three feasible request configurations and exposes a 6.04-point observation-success gain at load 20. The same change increases trajectory loss; at lower loads it does not improve observation success. An empty intersection for both marginal requirements is a meaningful engineering decision. Source replacement prevents correlated versions from acquiring the meaning of independent measurements.

Results separate prediction from knowledge predicates. Forecasting reachable outcomes improves loss relative to shallow computation on independent inputs, with less concentration in two exceptional samples. At ordinary loads, predicate removal leaves almost the same accuracy and flat scheduling wins nominal selection. Under additional workload, the time predicate changes actions and availability but increases longer-horizon exposure. The knowledge layer specifies and enforces the selected engineering requirement; it provides no universal accuracy margin.

Contract-based design, runtime assurance and binomial calibration are established. Their contribution here is the coupled localization interface: immutable-source versions, retained computation, receiver replay, event-constrained policy selection and refusal of inapplicable statements. The secondary catalogue analysis supports the usefulness of these distinctions, without constituting a new prospective trial or new statistical guarantee. The checker does not infer hidden shifts, guarantee worst-case execution time or convert a population bound into per-scan confidence. Stronger guarantees require additional sensing, timing characterization and assumptions.

Shared-kernel CART/RF controls isolate scheduling and publication; complete routes expose the competitiveness of Fixed2. Full-system comparisons, independent commercial-airframe geometry and physical sensor trials remain necessary for wider deployment claims. Embedded timing, communication, return variation, motion mismatch and field map alignment need measured validation before deployment cards can be issued.

Conditional deployment references in the supplement give assumed map-error and latency scenarios. They are calculations, rather than physical observations. A physical study should characterize these quantities, freeze its mixture and events, then calibrate the installation service actually used by the platform.
